\chapter{Outils numériques et synthèse}
\label{corriges:20}
\section*{Exercice 20.1 -- Tableur et recopie}
\texttt{B8 : =SOMME(B2:B6)}.
\texttt{C2 : =B2/\$B\$8}.
Moyenne : \texttt{=SOMMEPROD(A2:A6;B2:B6)/SOMME(B2:B6)}.
La référence \texttt{B2} devient \texttt{B3}, etc. en descendant. \texttt{\$B\$8} reste fixée : toutes les fréquences doivent être divisées par le même total.

\section*{Exercice 20.2 -- Construction dynamique}
Créer $A,B,C$ non alignés ; créer $M$ comme point du segment $[AB]$ ; tracer la parallèle à $(BC)$ par $M$ ; définir $N$ par intersection avec $(AC)$. Les dépendances imposent les alignements et le parallélisme. Les mesures suggèrent et vérifient numériquement l'égalité dans les cas testés ; elles ne prouvent pas tous les cas. La démonstration par les aires fournit l'argument général.

\section*{Exercice 20.3 -- Un modèle de remplissage}
La capacité est $\pi(0{,}5)^2\times2=0{,}5\pi\unite{m^3}=500\pi\unite L$. Le volume visé est $300\pi\unite L$, donc le temps est $15\pi\unite{min}\approx47{,}12$ minutes. Avant saturation, si le réservoir est vide au départ, $V(t)=20t$ litres pour $0\leq t\leq25\pi$ minutes. Au-delà, cette expression n'est plus le volume stocké.

\section*{Exercice 20.4 -- Aire maximale, complément}
L'autre côté vaut $20-x$, avec $0<x<20$. Donc $A(x)=x(20-x)=20x-x^2=100-(x^2-20x+100)=100-(x-10)^2$. Un carré étant positif ou nul, $A(x)\leq100$. L'égalité est atteinte pour $x=10$, donc le carré de côté $10\unite{cm}$ a l'aire maximale $100\unite{cm^2}$.
